% =============================================================================
% PREAMBLE.TEX
% LaTeX Preamble for Unit 1: Electrostatics - Class 12 Physics
% Board Exam Preparation Series
% Professional Design Guide Implementation - Modular Structure
% =============================================================================

% Note: \documentclass and \usetheme should be declared in main.tex, not here

% Typography: scientific sans-serif with crisp mathematics
\usefonttheme{professionalfonts}
\usepackage[sfdefault]{FiraSans}
\usepackage{newtxmath}
\renewcommand{\familydefault}{\sfdefault}

% Color Scheme - Unified and comprehensive
\definecolor{deepnavy}{HTML}{003366}        % Darker blue for structure
\definecolor{steelblue}{HTML}{4682B4}       % Main theme color
\definecolor{warmgray}{HTML}{696969}        % Dark gray for secondary text
\definecolor{lightblue}{HTML}{E0F0FF}       % Lighter blue for backgrounds

% Semantic Colors Based on User Syllabus Requirements
\definecolor{importantRed}{HTML}{D32F2F}    % Red for important/frequently asked
\definecolor{conceptualBlue}{HTML}{1976D2}  % Blue for conceptual understanding
\definecolor{familiarGreen}{HTML}{388E3C}   % Green to get familiar
\definecolor{commentPink}{HTML}{E91E63}     % Pink for comments/directions

% Legacy/Compatibility colors (unified)
\definecolor{primary}{HTML}{000080}         % Deep Navy (alias for deepnavy)
\definecolor{secondary}{HTML}{228B22}       % Forest Green
\definecolor{forestgreen}{HTML}{228B22}     % Forest Green (duplicate alias)
\definecolor{accentorange}{HTML}{FF4500}    % Orange for highlights
\definecolor{coolgray}{HTML}{555555}        % Neutral supporting text
\definecolor{slidebg}{HTML}{F5F5F5}         % Light slide backdrop
\definecolor{charcoal}{HTML}{333333}        % Dark gray for text

% Footer specific colors for four distinct sections
\definecolor{footerAuthor}{HTML}{1565C0}    % Strong blue for author section
\definecolor{footerChapter}{HTML}{2E7D32}   % Strong green for chapter section
\definecolor{footerPage}{HTML}{AD1457}      % Strong pink for page number
\definecolor{footerInstitute}{HTML}{E65100} % Strong orange for institute

% Beamer theme color settings
\setbeamercolor{background canvas}{bg=white}
\setbeamercolor{normal text}{fg=charcoal}
\setbeamercolor{frametitle}{bg=white, fg=steelblue!90!black}
\setbeamercolor{title}{fg=steelblue!90!black}
\setbeamercolor{structure}{fg=steelblue!80!black}
\setbeamercolor{alerted text}{fg=importantRed}
\setbeamercolor{block title}{fg=steelblue!80!black,bg=steelblue!12}
\setbeamercolor{block body}{fg=charcoal,bg=steelblue!5}
\setbeamercolor{titleinfobox}{fg=coolgray,bg=white}

% Professional footer colors
\setbeamercolor{footline}{fg=steelblue!90!black,bg=steelblue!10}
\setbeamercolor{author in head/foot}{fg=footerAuthor,bg=footerAuthor!10}
\setbeamercolor{chapter in head/foot}{fg=footerChapter,bg=footerChapter!10}
\setbeamercolor{page number in head/foot}{fg=footerPage,bg=footerPage!10}
\setbeamercolor{institute in head/foot}{fg=footerInstitute,bg=footerInstitute!10}

% Typography Settings - Following Design Guide with enhanced hierarchy
\setbeamerfont{title}{size=\Large, series=\bfseries}
\setbeamerfont{frametitle}{size=\large, series=\bfseries}
\setbeamerfont{normal text}{size=\small}
\setbeamerfont{author in head/foot}{size=\scriptsize,series=\bfseries}
\setbeamerfont{chapter in head/foot}{size=\scriptsize,series=\bfseries}
\setbeamerfont{page number in head/foot}{size=\scriptsize,series=\bfseries}
\setbeamerfont{institute in head/foot}{size=\scriptsize,series=\bfseries}

% Frame title template - Clean and scientific
\setbeamertemplate{frametitle}{%
    \nointerlineskip%
    \vspace{0.3cm}%
    {\color{steelblue!90}\Large\bfseries\insertframetitle}%
    \vspace{0.2cm}%
}

% Spacing Implementation - Optimized margins for 16:9 aspect ratio and better readability
% 16:9 aspect ratio: 25.6cm × 14.4cm (native dimensions for widescreen)
% Margins optimized for widescreen to prevent text crowding
\setbeamersize{text margin left=0.8cm, text margin right=0.8cm}

% 16:9 WIDESCREEN OPTIMIZATION SECTION
% ====================================
% For 16:9 presentations, optimal line length is 60-70 characters
% With 0.46\textwidth columns, each column provides perfect readability
% Vertical spacing adjusted to leverage extra horizontal space

% Enhanced vertical spacing for 16:9 (more horizontal space = less need for vertical)
\usepackage{setspace}
\setstretch{1.05}  % Slightly reduced from 1.08 for widescreen proportions

% Better paragraph spacing for widescreen
\parskip=0.4em plus 0.2em minus 0.2em

% Improved line spacing for mathematical content on widescreen
\jot=0.3em  % Reduced from default for tighter math displays

% Add slide numbers and navigation symbols
\setbeamertemplate{navigation symbols}{}
\setbeamerfont{page number in head/foot}{size=\scriptsize}

% Enhanced packages
\usepackage{tikz}
\usepackage{amsmath}
\usepackage{mathtools}
\usetikzlibrary{arrows.meta,positioning,decorations.pathreplacing,calc,shapes.geometric,shadows,decorations.markings}
\usepackage{graphicx}
\usepackage{circuitikz}
\usepackage{pgfplots}
\pgfplotsset{compat=1.18}
\usepackage[most]{tcolorbox}
\usepackage{fontawesome5}
\usepackage{booktabs}
\usepackage{array}
\usepackage{multirow}
\usepackage{tabularx}
\usepackage{amsmath,amssymb,amsthm}
\usepackage{siunitx}
\usepackage{physics}
\sisetup{inter-unit-product=\ensuremath{{}\cdot{}}}

% Fix for siunitx and physics package conflict - ensure \qty is properly defined
\AtBeginDocument{\RenewCommandCopy\qty\SI}

% Enhanced visual commands for educational content
% Custom commands for semantic highlighting (unified definitions)
\newcommand{\important}[1]{\textcolor{importantRed}{\textbf{#1}}}   % Red for important/frequently asked
\newcommand{\concept}[1]{\textcolor{conceptualBlue}{\textbf{#1}}}   % Blue for concepts
\newcommand{\familiar}[1]{\textcolor{familiarGreen}{#1}}             % Green to get familiar
\newcommand{\directed}[1]{\textcolor{commentPink}{\textit{#1}}}     % Pink for comments/directions
\newcommand{\formula}[1]{\colorbox{lightblue!40}{$\displaystyle #1$}} % Formula highlighting
\newcommand{\examplenew}[1]{\textcolor{steelblue}{\textbf{#1}}}     % Example highlighting
% Backward-compatible alias (define only if not already defined by a package)
\providecommand{\example}[1]{\examplenew{#1}}
\newcommand{\keypoint}[1]{\textcolor{familiarGreen}{\textbf{#1}}}   % Key points

% Chapter name command for footer
\newcommand{\currentchapter}{Chapter Name}  % Default value
\newcommand{\setchapter}[1]{\renewcommand{\currentchapter}{#1}}

% Enhanced section titles
\newcommand{\sectiontitle}[2]{%
    \begin{tcolorbox}[colback=deepnavy, colframe=deepnavy,
        boxrule=0pt, arc=3mm, boxsep=4pt,
        left=8pt, right=8pt, top=4pt, bottom=4pt,
        valign=center]
        \textcolor{white}{\Large\textbf{#1: #2}}%
    \end{tcolorbox}%
}

% Board exam focus box
\newtcolorbox{exambox}[1]{
    colback=yellow!20,
    colframe=accentorange,
    colbacktitle=accentorange!80,
    coltitle=white,
    fonttitle=\bfseries,
    title=#1,
    boxrule=2pt,
    arc=2mm,
    boxsep=6pt,
    left=8pt,
    right=8pt,
    top=4pt,
    bottom=4pt
}

% Professional content boxes - Enhanced scientific styling
\newtcolorbox{keybox}[1][Key Point]{
    colback=white,
    colframe=steelblue!85,
    coltitle=white,
    colbacktitle=steelblue!85,
    fonttitle=\bfseries\normalsize,
    title=#1,
    boxrule=2pt,
    arc=3mm,
    boxsep=8pt,
    left=10pt,
    right=10pt,
    top=6pt,
    bottom=6pt,
    fontupper=\normalsize,
    enhanced,
    shadow={2pt}{2pt}{0pt}{black!20}
}

\newtcolorbox{importantbox}[1][Important]{
    colback=white,
    colframe=importantRed!85,
    coltitle=white,
    colbacktitle=importantRed!85,
    fonttitle=\bfseries\normalsize,
    title=#1,
    boxrule=2pt,
    arc=3mm,
    boxsep=8pt,
    left=10pt,
    right=10pt,
    top=6pt,
    bottom=6pt,
    fontupper=\normalsize,
    enhanced,
    shadow={2pt}{2pt}{0pt}{black!15}
}

\newtcolorbox{examplebox}[1][Example]{
    colback=white,
    colframe=familiarGreen!80,
    coltitle=white,
    colbacktitle=familiarGreen!80,
    fonttitle=\bfseries\normalsize,
    title=#1,
    boxrule=2pt,
    arc=3mm,
    boxsep=8pt,
    left=10pt,
    right=10pt,
    top=6pt,
    bottom=6pt,
    fontupper=\normalsize,
    enhanced,
    shadow={2pt}{2pt}{0pt}{black!15}
}

\newtcolorbox{formulabox}[1][Formula]{
    colback=steelblue!8,
    colframe=steelblue!90,
    coltitle=white,
    colbacktitle=steelblue!90,
    fonttitle=\bfseries\normalsize,
    title=#1,
    boxrule=2pt,
    arc=3mm,
    boxsep=8pt,
    left=10pt,
    right=10pt,
    top=6pt,
    bottom=6pt,
    fontupper=\large\centering,
    enhanced,
    shadow={2pt}{2pt}{0pt}{black!20}
}

% Two-column layout support for overflow content - Optimized for 16:9
% Each column: 0.46\textwidth + 0.08\textwidth gap = 1.0\textwidth
\newcommand{\twocolumncontent}[2]{%
    \begin{columns}[T]
        \column{0.46\textwidth}
        #1
        \column{0.46\textwidth}
        #2
    \end{columns}
}

% Three-column layout support - Optimized for 16:9 widescreen
% Each column: 0.29\textwidth + gaps = proper widescreen distribution
\newcommand{\threecolumncontent}[3]{%
    \begin{columns}[T]
        \column{0.29\textwidth}
        #1
        \column{0.29\textwidth}
        #2
        \column{0.29\textwidth}
        #3
    \end{columns}
}

% Custom macros for specific chapters (can be adapted)
\newcommand{\ZT}{\ensuremath{ZT}}
\newcommand{\Seebeck}{\ensuremath{\alpha}}
\newcommand{\conductivity}{\ensuremath{\sigma}}
\newcommand{\resistivity}{\ensuremath{\rho}}
\newcommand{\kappaT}{\ensuremath{\kappa}}

% Custom FrameColorBox environment for compatibility
\newenvironment{FrameColorBox}[2]{%
    \begin{beamercolorbox}[bg=#1,#2]{block body}%
}{%
    \end{beamercolorbox}%
}

% Graphics path and extensions
\graphicspath{{./}{Images/}{images/}}
\DeclareGraphicsExtensions{.png,.jpg,.jpeg,.pdf}

% Professional logo commands - Using actual logo path
\newcommand{\presentationlogo}{%
    \IfFileExists{Images/logo.png}{%
        \includegraphics[height=0.5cm]{Images/logo.png}%
    }{%
        \raisebox{0pt}[0pt][0pt]{\scalebox{1}{\tikz[baseline=(current bounding box.center)]{\draw[fill=white,draw=steelblue,line width=1.1pt] (0,0) circle (0.18cm);\node at (0,0) [font=\tiny\bfseries,steelblue] {RP};}}}%
    }%
}

% Small logo for footer
\newcommand{\footerlogo}{%
    \IfFileExists{Images/logo.png}{%
        \includegraphics[height=0.25cm]{Images/logo.png}%
    }{%
        \raisebox{0pt}[0pt][0pt]{\scalebox{1}{\tikz[baseline=(current bounding box.center)]{\draw[fill=white,draw=steelblue,line width=1.1pt] (0,0) circle (0.12cm);\node at (0,0) [font=\tiny\bfseries,steelblue] {RP};}}}%
    }%
}

% Bigger logo for title slide
\newcommand{\biglogo}{%
    \IfFileExists{Images/logo.png}{%
        \includegraphics[height=0.8cm]{Images/logo.png}%
    }{%
        \fcolorbox{steelblue}{white}{\textcolor{steelblue}{\textbf{RP}}}%
    }%
}

% Fallback logo command in case the image doesn't exist
\newcommand{\fallbacklogo}{%
    \tikz[baseline=(current bounding box.center)]{
        \draw[fill=white,draw=steelblue,line width=1.1pt] (0,0) circle (0.38cm);
        \node at (0,0) [font=\tiny\bfseries,steelblue] {RP};
    }
}

% Small electron animation for footer
\newcommand{\electronanimation}{%
    \begin{tikzpicture}[scale=0.3]
        % Nucleus
        \fill[red!80] (0,0) circle (0.15);
        % Electron orbit
        \draw[gray!50] (0,0) circle (0.6);
        % Animated electron
        \node[circle,fill=blue,inner sep=1.5pt] (e) at (0.6,0) {};
        % Electron path arrow
        \draw[->,blue!50,decoration={markings,mark=at position .25 with {\arrow{>}}},postaction={decorate}] (0.6,0) arc (0:270:0.6);
        % Label
        \node[font=\tiny,gray] at (0,-0.9) {e-};
    \end{tikzpicture}%
}

% Minimal frametitle styling with underline
\setbeamertemplate{frametitle}{%
    \vspace*{0.5em}%
    \begin{beamercolorbox}[wd=\paperwidth,ht=3.2ex,dp=1.1ex,leftskip=1.1em,rightskip=1.1em]{frametitle}
        \usebeamerfont{frametitle}\insertframetitle
    \end{beamercolorbox}
    \nointerlineskip
    {\color{steelblue!70}\rule{\paperwidth}{0.6pt}}\vspace{-0.35em}
}

% Helper environment for schematic placeholders
\newenvironment{conceptdiagram}[2][0.85\textwidth]{%
    \begin{center}
    \begin{tcolorbox}[width=#1,colback=steelblue!7,colframe=steelblue!70,
        boxrule=0.9pt,arc=1.6mm,left=8pt,right=8pt,top=6pt,bottom=6pt,
        title=\textbf{#2}]
    \small
}{%
    \end{tcolorbox}
    \end{center}
}

% 16:9 WIDESCREEN SPECIFIC ENVIRONMENTS
% =====================================
% For better content management in widescreen format

% Optimized single-column frame content (full width)
\newenvironment{fullwidthcontent}{%
    \begin{center}
    \begin{minipage}{0.92\textwidth}
}{%
    \end{minipage}
    \end{center}
}

% Show a small logo on every slide using beamer's \logo template
\logo{\presentationlogo}

% Header: show a small logo at the top-left for normal slides (hidden on plain/title frames)
\makeatletter
\setbeamertemplate{headline}{%
    \ifx\insertframetitle\@empty
        % plain frame or no frametitle: don't show headline
    \else
        \begin{beamercolorbox}[wd=\paperwidth,sep=0pt,ht=1.2cm,dp=0.2cm]{palette quaternary}%
            % Header without logo
        \end{beamercolorbox}
    \fi
}
\makeatother

% Enhanced footer with four distinct sections, larger font, logo, and electron animation
\setbeamertemplate{footline}{%
    \leavevmode
    \hbox{%
        % Author section with logo
        \begin{beamercolorbox}[wd=.24\paperwidth,ht=4.5ex,dp=1.5ex,left,leftskip=0.2cm]{author in head/foot}%
            \usebeamerfont{author in head/foot}\footerlogo\hspace{0.1cm}\textbf{Dr. KN Ray}
        \end{beamercolorbox}%
        % Chapter section
        \begin{beamercolorbox}[wd=.24\paperwidth,ht=4.5ex,dp=1.5ex,center]{chapter in head/foot}%
            \usebeamerfont{chapter in head/foot}\textbf{\insertshorttitle}
        \end{beamercolorbox}%
        % Page number section with electron animation
        \begin{beamercolorbox}[wd=.24\paperwidth,ht=4.5ex,dp=1.5ex,center]{page number in head/foot}%
            \usebeamerfont{page number in head/foot}\electronanimation\hspace{0.1cm}\textbf{\insertframenumber{} / \inserttotalframenumber}
        \end{beamercolorbox}%
        % Institute section
        \begin{beamercolorbox}[wd=.24\paperwidth,ht=4.5ex,dp=1.5ex,right,rightskip=0.2cm]{institute in head/foot}%
            \usebeamerfont{institute in head/foot}\textbf{Ray Physics Academy}
        \end{beamercolorbox}%
    }%
    \vskip0pt%
}

% =============================================================================
% CUSTOM FOOTER CONFIGURATION
% =============================================================================

% Define lighter golden and yellow colors for footer sections
\definecolor{goldenhue}{HTML}{FFE680}      % Lighter golden color for author section
\definecolor{brightyellow}{HTML}{FFF59D}   % Lighter yellow color for chapter section

% Define footer template with 4 sections - enhanced styling
\setbeamertemplate{footline}{
    \leavevmode%
    \hbox{%
        % Section 1: Institute name with logo (27% width) - Keep as-is (PERFECT)
        \begin{beamercolorbox}[wd=.27\paperwidth,ht=4ex,dp=1.5ex,left,leftskip=0.5em]{author in head/foot}%
            \usebeamerfont{author in head/foot}%
            \IfFileExists{Images/logo.png}{%
                \raisebox{-0.6ex}{\includegraphics[height=3ex]{Images/logo.png}}%
            }{%
                \raisebox{-0.6ex}{\tikz[baseline=(current bounding box.center)]{\draw[fill=white,draw=steelblue,line width=1pt] (0,0) circle (0.15cm);\node at (0,0) [font=\tiny\bfseries,steelblue] {RP};}}%
            }%
            \hspace{0.4em}%
            \textcolor{steelblue}{\small\textbf{Ray Physics Academy}}
        \end{beamercolorbox}%
        %
        % Section 2: Author name with photo and LIGHTER GOLDEN background (27% width)
        \begin{beamercolorbox}[wd=.27\paperwidth,ht=4ex,dp=1.5ex,center]{title in head/foot}%
            \colorbox{goldenhue}{%
                \parbox[c][4ex][c]{0.265\paperwidth}{%
                    \centering
                    \usebeamerfont{title in head/foot}%
                    \IfFileExists{Images/knray.png}{%
                        \raisebox{-0.6ex}{\includegraphics[height=3ex,width=3ex,keepaspectratio]{Images/knray.png}}%
                        \hspace{0.3em}%
                    }{}%
                    \textcolor{black}{\small\textbf{Dr. KN Ray}}%
                }%
            }%
        \end{beamercolorbox}%
        %
        % Section 3: Chapter name with LIGHTER YELLOW background (32% width)
        \begin{beamercolorbox}[wd=.32\paperwidth,ht=4ex,dp=1.5ex,center]{date in head/foot}%
            \colorbox{brightyellow}{%
                \parbox[c][4ex][c]{0.315\paperwidth}{%
                    \centering
                    \usebeamerfont{date in head/foot}%
                    \textcolor{black}{\small\textbf{\currentchapter}}%
                }%
            }%
        \end{beamercolorbox}%
        %
        % Section 4: Slide number (14% width) - Keep as-is (PERFECT)
        \begin{beamercolorbox}[wd=.14\paperwidth,ht=4ex,dp=1.5ex,right,rightskip=0.5em]{page number in head/foot}%
            \usebeamerfont{page number in head/foot}%
            \textcolor{charcoal}{\small\textbf{Slide \insertframenumber{} / \inserttotalframenumber}}%
        \end{beamercolorbox}%
    }%
    \vskip0pt%
}

% Remove navigation symbols
\setbeamertemplate{navigation symbols}{}

% =============================================================================

% Additional packages
\usepackage{textcomp}
\usepackage[absolute,overlay]{textpos}
\usepackage{caption}
\usepackage{subcaption}
\usepackage{float}
\usepackage{lipsum}
\usepackage{fancyvrb}
\usepackage{listings}
\usepackage{color}
\usepackage{colortbl}
\usepackage{longtable}
\usepackage{appendixnumberbeamer}

\usepackage{setspace}
\usepackage{hyperref}
\usepackage{wallpaper}
\usepackage{bookmark}

% Declare Unicode characters for subscripts
\DeclareUnicodeCharacter{2082}{\textsubscript{2}}
\DeclareUnicodeCharacter{2083}{\textsubscript{3}}

% Note: \setstretch already defined above in 16:9 optimization section
% Value: 1.05 (optimized for widescreen, reduced from default 1.08)

\hypersetup{
    colorlinks=true,
    linkcolor=deepnavy,
    citecolor=forestgreen,
    urlcolor=forestgreen,
    pdfpagemode=UseOutlines,
    pdfstartview=FitH,
    pdftitle={Unit 1: Electrostatics - Class 12 Physics},
    pdfauthor={Dr. KN Ray},
    pdfsubject={Board Exam Preparation Notes},
    pdfkeywords={Electrostatics, Physics, Class 12, CBSE, ICSE}
}

% Enhanced list formatting per Design Guide
\setbeamertemplate{itemize item}{\textcolor{deepnavy}{$\bullet$}}
\setbeamertemplate{itemize subitem}{\textcolor{steelblue}{--}}
\setbeamertemplate{enumerate item}{\textcolor{deepnavy}{\arabic{enumi}.}}
\setbeamertemplate{description item}{\textcolor{deepnavy}{\textbf{\insertdescriptionitem}}}

% Background image configuration for 16:9 slides
% \CenterWallPaper{1}{Images/background.png}

% Enhanced caption formatting
\captionsetup{
    font=small,
    labelformat=empty,
    labelsep=none,
    textfont=it,
    justification=centering,
    singlelinecheck=false
}